\clearpage
\fancyhead[L]{\small GAUSS AN BODE.}
\SourceLeaf{456}

\begin{center}
{\large GAUSS AN BODE.}
\end{center}

\begin{center}
\rule{6cm}{0.4pt}\\
Astronomisches Jahrbuch für 1828. Seite 171..172. Berlin 1825.\\
\rule{6cm}{0.4pt}
\end{center}

\begin{flushright}
Göttingen, den 5. August 1825.
\end{flushright}

Diese Beobachtungen sind vielleicht deshalb von größerem Interesse, weil diese Planeten 1825 zum ersten Mal wieder ganz in der Nähe der Gegend standen, in der sie 1802 entdeckt beziehungsweise wiedergefunden wurden. Zumindest Pallas wurde vielleicht auch sonst nirgends beobachtet, da die veröffentlichte Ephemeride so stark abweicht. Mir selbst hatte diesmal die Zeit gefehlt, eine solche nach meiner Theorie zu berechnen.

\begin{center}
\emph{Meridianbeobachtungen der Pallas an der Sternwarte in Göttingen.}
\end{center}

\begin{center}
\begingroup\small\setlength{\tabcolsep}{4pt}
\begin{tabular}{lr|c|c|c}
\multicolumn{2}{c|}{1825} & \shortstack{Mittlere Zeit\\in Göttingen} & Rektaszension & Deklination \\
\hline
März & 3 & $13^{\rm u}\,4^{\rm m}\,23^{\rm s}7$ & $177^{\circ}\,34'\,24''3$ & $0^{\circ}\,24'\,4''1$ Südl. \\
 & 4 & $12\quad59\quad49.6$ & $177\quad24\quad50.7$ & $0\quad0\quad18.2$ Nordl. \\
 & 7 & $12\quad46\quad4.0$ & $176\quad55\quad18.1$ & $1\quad14\quad7.2$ \\
 & 8 & $12\quad41\quad27.9$ & $176\quad45\quad12.3$ & $1\quad38\quad52.6$ \\
 & 9 & $12\quad36\quad51.3$ & $176\quad35\quad0.3$ & $2\quad3\quad40.3$ \\
 & 10 & $12\quad32\quad14.3$ & $176\quad24\quad43.2$ & $2\quad28\quad29.3$ \\
 & 18 & $11\quad55\quad14.1$ & $175\quad1\quad15.0$ & $5\quad44\quad59.3$ \\
 & 19 & $11\quad50\quad37.3$ & $174\quad50\quad59.1$ & $6\quad8\quad56.3$ \\
 & 20 & $11\quad46\quad0.6$ & $174\quad40\quad46.2$ & $6\quad32\quad41.4$ \\
 & 25 & $11\quad23\quad6.1$ & $173\quad51\quad53.7$ & $8\quad27\quad50.6$ \\
 & 27 & $11\quad14\quad1.1$ & $173\quad33\quad32.8$ & $9\quad11\quad51.7$ \\
 & 29 & $11\quad4\quad59.8$ & $173\quad16\quad8.2$ & $9\quad54\quad29.8$ \\
April & 6 & $10\quad29\quad36.6$ & $172\quad16\quad59.2$ & $12\quad29\quad56.1$ \\
 & 7 & $10\quad25\quad16.7$ & $172\quad10\quad57.4$ & $12\quad47\quad32.9$ \\
\end{tabular}
\endgroup
\end{center}

\begin{center}
\emph{Meridianbeobachtungen der Ceres an der Sternwarte in Göttingen.}
\end{center}

\begin{center}
\begingroup\small\setlength{\tabcolsep}{4pt}
\begin{tabular}{lr|c|c|c}
\multicolumn{2}{c|}{1825} & \shortstack{Mittlere Zeit\\in Göttingen} & Rektaszension & Deklination \\
\hline
März & 9 & $13^{\rm u}\,1^{\rm m}\,7^{\rm s}5$ & $182^{\circ}\,40'\,4''6$ & $17^{\circ}\,35'\,18''5$ Nordl. \\
 & 10 & $12\quad56\quad23.3$ & $182\quad27\quad57.4$ & $17\quad41\quad20.9$ \\
 & 18 & $12\quad18\quad11.1$ & $180\quad46\quad26.4$ & $18\quad24\quad4.4$ \\
 & 19 & $12\quad13\quad23.5$ & $180\quad33\quad28.9$ & $18\quad28\quad34.0$ \\
 & 20 & $12\quad8\quad36.8$ & $180\quad20\quad32.2$ & $18\quad32\quad50.0$ \\
 & 29 & $11\quad25\quad38.4$ & $178\quad26\quad37.3$ & $19\quad1\quad18.0$ \\
April & 7 & $10\quad43\quad31.3$ & $176\quad45\quad21.7$ & $19\quad10\quad5.9$ \\
\end{tabular}
\endgroup
\end{center}
